% Composition g-modes of neutron stars: theory, observability, computation.
% The paper-style description of eos/astro/gmode (CLAUDE.md section 11).
% Compile with:  pdflatex gmode && bibtex gmode && pdflatex gmode
\documentclass[aps,prd,twocolumn,superscriptaddress,nofootinbib,10pt]{revtex4-2}

\usepackage{amsmath,amssymb}
\usepackage{graphicx}
\usepackage[colorlinks=true,allcolors=blue]{hyperref}

\begin{document}

\title{Composition $g$-modes as a probe of dense-matter composition}

\author{}
\affiliation{}
\date{\today}

\begin{abstract}
The gravity ($g$) modes of a neutron star are restored by buoyancy, which
exists only when a displaced fluid element cannot chemically re-equilibrate
within an oscillation period. Their frequencies therefore measure the
\emph{composition gradient} of dense matter, an aspect of the equation of state
to which mass, radius and tidal deformability are blind. We summarise the
theory, the conditions under which such a mode could be detected, and its
numerical computation, and we examine in detail which reactions are fast enough
to matter in a quark--hadron mixed phase --- equivalently, which quantities must
be held fixed in the thermodynamic derivatives that define the two sound speeds.
\end{abstract}

\maketitle

\section{Buoyancy and the two sound speeds}

Displace a fluid element radially outwards by $\xi_r$ in a stratified star. It
re-establishes pressure balance with its new surroundings essentially
instantaneously, on the sound-crossing time, but its \emph{composition} relaxes
only as fast as the reactions that change particle fractions. If those
reactions are slow, the element arrives with the composition it had, which is
in general not the equilibrium composition at its new location. It is therefore
denser or lighter than its surroundings, and the resulting buoyancy force
drives an oscillation~\cite{ReiseneggerGoldreich1992}.

The strength of this restoring force is the difference between two distinct
sound speeds. The \emph{equilibrium} sound speed
\begin{equation}
  c_{\rm eq}^{2} = \frac{dP}{d\varepsilon}
\end{equation}
is the total derivative along the equilibrated sequence, with the composition
free to track the density; this is the quantity that enters the
Tolman--Oppenheimer--Volkoff equations and that a mass--radius measurement
constrains. The \emph{frozen}, or adiabatic, sound speed
\begin{equation}
  c_{\rm ad}^{2} = \left(\frac{\partial P}{\partial \varepsilon}\right)_{\!x}
\end{equation}
is taken at fixed composition $x$. The Brunt--V\"ais\"al\"a frequency built from
them,
\begin{equation}
  N^{2} = g^{2}\left(\frac{1}{c_{\rm eq}^{2}}
                   - \frac{1}{c_{\rm ad}^{2}}\right) e^{\nu-\lambda},
  \qquad g = -\frac{1}{\varepsilon+P}\frac{dP}{dr},
  \label{eq:bv}
\end{equation}
is the local frequency of that oscillation~\cite{Jaikumar2021}. Here $\nu$ and
$\lambda$ are the metric functions of the static interior,
$ds^{2}=-e^{\nu}dt^{2}+e^{\lambda}dr^{2}+r^{2}d\Omega^{2}$, with
$e^{\lambda}=(1-2Gm/rc^{2})^{-1}$.

Equation~(\ref{eq:bv}) is the entire point. If matter has a single sound speed
--- a one-component fluid, or any model without a density-dependent composition
--- then $N^{2}\equiv0$ and \emph{no composition $g$-mode exists}. Convective
stability requires $c_{\rm ad}>c_{\rm eq}$, and typically
$c_{\rm ad}^{2}-c_{\rm eq}^{2}\sim10^{-3}$--$10^{-1}$, so $g$-mode frequencies
lie far below the $\sim2$~kHz fundamental ($f$) mode, at
$\sim100$--$300$~Hz for nucleonic matter.

The diagnostic power comes from what happens at a phase transition. Across a
quark--hadron mixed phase the appearance of a new species lets the system absorb
compression by converting hadrons into quarks rather than by raising the
pressure, so $c_{\rm eq}$ \emph{drops sharply}, while $c_{\rm ad}$, evaluated at
frozen composition, does not. The difference in Eq.~(\ref{eq:bv}) jumps, and
with it the $g$-mode frequency, which for a hybrid star can exceed $500$~Hz ---
a factor of two or more above a purely nucleonic star of the same
mass~\cite{Jaikumar2021,Constantinou2021}. Since hybrid and nucleonic stars can
be constructed with indistinguishable $M$, $R$ and $\Lambda$, this is
information no static observable supplies.

\section{Which reactions are frozen}

The frozen sound speed is not defined until one states what is held fixed, and
that follows from comparing reaction timescales with the oscillation period, a
few milliseconds for a $g$-mode. Writing $\tau$ for the equilibration time and
$\omega$ for the mode frequency, a channel is \emph{frozen} when
$\omega\tau\gg1$ and \emph{equilibrated} when $\omega\tau\ll1$.

\paragraph{Strong and electromagnetic processes ($\tau\sim10^{-23}$~s).}
Always fast. In a two-phase mixture this is what keeps the phases in
\emph{mechanical} equilibrium at a common pressure throughout the oscillation,
a fact we use below. Flavour-conserving rearrangements within each phase also
equilibrate.

\paragraph{Non-leptonic weak processes in quark matter ($\tau\sim10^{-8}$~s).}
The reaction $u+d\leftrightarrow u+s$ is fast on a millisecond
period~\cite{Madsen1993}. Strangeness in the quark phase therefore
\emph{re-equilibrates} and arguably should not be frozen --- a point on which
implementations differ, and to which we return.

\paragraph{Leptonic weak processes (Urca, $\tau\propto T^{-4}$ or $T^{-6}$).}
Direct Urca $n\to p+e+\bar\nu_{e}$ and its modified counterpart are slow for
$T\lesssim1$~MeV. \emph{Their slowness is why composition $g$-modes exist at
all} in cold neutron stars: the lepton fractions, and hence the proton fraction,
are frozen.

\paragraph{Quark--hadron conversion ($\tau\gg$~ms).}
Changing the quark volume fraction $\chi$ requires transferring baryon number
across phase boundaries with an accompanying change of strangeness --- a
high-order weak process, further limited by transport across the mixed-phase
structures. It is slow, so \textbf{$\chi$ is frozen}. This is the load-bearing
assumption behind the large mixed-phase $g$-mode: were $\chi$ free to relax,
$c_{\rm ad}$ would collapse onto $c_{\rm eq}$ and the signal would vanish.

In short, the derivative defining $c_{\rm ad}^{2}$ is taken at fixed $\chi$ and
fixed lepton fractions, with pressure equality between phases and the internal
strong-interaction composition of each phase left free.

\subsection{Finite rates}
\label{sec:rates}

The frozen/equilibrated dichotomy is a limit of a continuous problem. For a
finite equilibration rate $\gamma$ the response is neither, but a complex,
frequency-dependent interpolation~\cite{Counsell2025}
\begin{equation}
  c_{\rm dy}^{2} = c_{\rm eq}^{2}
    + \frac{c_{\rm ad}^{2}-c_{\rm eq}^{2}}{1+\gamma/i\omega},
  \label{eq:cdy}
\end{equation}
whose real part interpolates between the two limits and whose imaginary part,
maximal at $\gamma=\omega$, is bulk viscosity,
$\zeta=(\varepsilon+P)\,{\rm Im}\,c_{\rm dy}^{2}/\omega$. Substituting
Eq.~(\ref{eq:cdy}) into Eq.~(\ref{eq:bv}) makes $N^{2}$ and hence the
eigenfrequency complex, the imaginary part being the bulk-viscous damping rate.
The suppression is progressive. For a $1.4\,M_{\odot}$ DD2~\cite{Typel2010}
star whose undamped fundamental $g$-mode sits at $150$~Hz, we obtain a quality
factor $Q=\mathrm{Re}\,\omega/2\,\mathrm{Im}\,\omega$ of $10.4$, $3.4$, $0.96$
and $0.53$ at $\gamma/\omega=0.1$, $0.3$, $1$ and $3$, the mode frequency
falling from $149$ to $27$~Hz as the buoyancy is eaten away; by
$\gamma\simeq10\,\omega$ no mode remains. The crossing into critical damping,
$Q\simeq1$, coincides with $\gamma\simeq\omega$ as expected. Using
Fermi-surface Urca rates~\cite{AlfordHarris2018,Alford2020} we find
$\gamma=\omega$ at $T\simeq5$~MeV for a 1~kHz oscillation, nearly independent
of density, consistent with~\cite{AlfordHarris2019}. Cold stars, $T\ll1$~MeV,
are firmly frozen, with $\gamma$ many orders of magnitude below $\omega$.

\section{Convention dependence: a caveat}

Two choices in the definition of the sound speeds change the answer materially,
and neither is settled.

\emph{(i) The lepton content.} Both sound speeds must describe the \emph{same}
fluid. If $c_{\rm eq}^{2}$ is computed along the $\beta$-equilibrium sequence of
charge-neutral matter (leptons included) while $c_{\rm ad}^{2}$ is computed for
the baryons alone, the difference is contaminated by the lepton contribution,
which is comparable to the entire buoyancy signal. In our DD2 calculation the
mismatched pairing yields $c_{\rm ad}^{2}<c_{\rm eq}^{2}$ over part of the
density range --- a spurious convective instability, $N^{2}<0$ --- whereas the
consistent pairing gives $N^{2}>0$ throughout. Consistency here is not a
refinement but a correctness condition.

\emph{(ii) Combining the phases in a mixed phase.} Because strong interactions
hold the two phases at a common pressure, a compression imposes the same $\delta
P$ on both, and each responds with its own $\delta\varepsilon$. With
$\varepsilon_{\rm mix}=(1-\chi)\varepsilon_{H}+\chi\varepsilon_{Q}$ at
$P_{H}=P_{Q}$ this gives the \emph{reciprocal}, volume-fraction-weighted
combination
\begin{equation}
  \frac{1}{c_{\rm ad,mix}^{2}} = \frac{1-\chi}{c_{\rm ad,H}^{2}}
                               + \frac{\chi}{c_{\rm ad,Q}^{2}},
  \label{eq:isobaric}
\end{equation}
pressure adding across the phases like a voltage across parallel
resistors~\cite{Jaikumar2021}. An alternative prescription compresses both
phases by the same density factor at fixed $\chi$ and volume-averages the
pressures; it is a defensible definition of ``frozen'', but it does \emph{not}
maintain a common pressure under the perturbation and so is inconsistent with
the mechanical equilibrium that the mode equations assume. The two agree in the
pure phases and differ only inside the mixed-phase window --- precisely where
the signal lives. A related open question is whether the quark phase's
strangeness should be frozen at all, given the fast non-leptonic rate noted
above; freezing it overestimates the buoyancy, and a fully
reaction-constrained treatment freezes only those species whose gradients
genuinely survive an oscillation period~\cite{Haskell2024}. We therefore regard
the mixed-phase $g$-mode frequency as carrying a systematic uncertainty of this
origin, and recommend quoting results under both prescriptions.

\section{What the equation of state must supply}
\label{sec:contract}

Equations~(\ref{eq:isobaric}) and the buoyancy above require exactly two
numbers at every point of the stellar sequence: $c_{\rm eq}^{2}$ and
$c_{\rm ad}^{2}$. Nothing else about the composition enters. In particular no
derivative $dY_{i}/dn_{B}$ is needed, because the composition gradient has
already been integrated into the difference of the two speeds --- which is why
the $g$-mode is often described as a ``tale of two sounds''~\cite{Jaikumar2021}.

The implementation makes that the literal interface. A model or a composite
engine produces the two columns alongside the $(P,\varepsilon,n_{B})$ table a
structure solver integrates, and the mode solver consumes them; the mode solver
imports no model, and no model imports the mode solver. Both columns come from
the model's own second-derivative entry point, one at fixed composition and one
along the equilibrated sequence.

The tables are at $T=0$. The condition that defines the frozen speed is
``without varying chemical composition'', not the vanishing of the temperature:
at $T=0$ along a sequence whose composition changes with density the two speeds
still differ, and $N^{2}\neq0$. Zero temperature removes only the \emph{thermal}
axis --- the distinction between an isothermal and an adiabatic second
derivative, which differ by $C_{P}/C_{V}$ --- leaving the composition axis, on
which the mode lives, entirely intact. A point therefore carries exactly two
numbers, and neither requires a label saying which thermal variable was held.
Finite temperature reintroduces that axis and is left to future work; it is also
the regime in which the finite-rate treatment of Sec.~\ref{sec:rates} matters,
the two being aspects of the same physics.

\section{Numerical computation}

We work in the relativistic Cowling approximation, holding the spacetime fixed
and perturbing only the fluid~\cite{McDermott1983}. This reduces the
fourth-order system of Thorne and Campolattaro~\cite{ThorneCampolattaro1967} to
two first-order equations confined to the star, at the cost of the
gravitational-wave damping time; for $g$-modes the frequency error is a few per
cent, since the mode is concentrated in the core and nearly divergence-free.

With $e^{-i\omega t}$ time dependence, spherical-harmonic degree $l$, and
variables $U=r^{2}e^{\lambda/2}\xi_{r}$ and $V=\delta P/(\varepsilon+P)$, the
equations are~\cite{Jaikumar2021}
\begin{align}
  \frac{dU}{dr} &= \frac{g}{c_{\rm ad}^{2}}U
    + e^{\lambda/2}\left[\frac{l(l+1)e^{\nu}}{\omega^{2}}
      - \frac{r^{2}}{c_{\rm ad}^{2}}\right]V, \nonumber\\
  \frac{dV}{dr} &= e^{\lambda/2-\nu}\frac{\omega^{2}-N^{2}}{r^{2}}U
    + g\left(\frac{1}{c_{\rm eq}^{2}}-\frac{1}{c_{\rm ad}^{2}}\right)V.
  \label{eq:cowling}
\end{align}
Note which speed appears where: $c_{\rm ad}$ governs the element's response to
compression, while $c_{\rm eq}$ enters only through the buoyancy terms.

The background is obtained by integrating the TOV equations while retaining the
radial profiles. The metric function $\nu$ needs no equation of its own, since
hydrostatic equilibrium gives $d\nu/dr=-2P'/(\varepsilon+P)$; it follows by
quadrature from the pressure profile, fixed by matching $e^{\nu(R)}=1-2GM/Rc^{2}$
at the surface, and the same relation supplies $g=\nu'/2$.

Regularity at the centre gives $U\to r^{l+1}$ and
$V\to(\omega^{2}/l)r^{l}e^{-\nu(0)}$, and the outer boundary condition is the
vanishing of the Lagrangian pressure perturbation,
$\Delta P=\delta P+\xi_{r}\,dP/dr=0$, i.e.
\begin{equation}
  V(R) = g(R)\,U(R)\,e^{-\lambda(R)/2}/R^{2}.
\end{equation}
Integrating outwards from the regular series always succeeds; the surface
condition holds only for discrete $\omega$, so one scans $\omega$, brackets sign
changes of the surface residual and refines them. The crust, which carries no
tabulated composition, is treated as a homogeneous fluid with
$c_{\rm ad}=c_{\rm eq}$ and hence $N^{2}=0$.

Modes are classified by the node count of $\xi_{r}$: none for the $f$-mode, $k$
for $g_{k}$ and $p_{k}$, with the $g$-branch below the $f$-mode and the
$p$-branch above. In practice it is more robust to separate the branches by the
buoyancy scale, since the $g$-branch is confined to $\omega<\max N$ while the
$f$-mode sits near the dynamical frequency $\sqrt{GM/R^{3}}$; node counting near
the surface is sensitive to where the star is truncated. For a complex $N^{2}$
the eigenfrequency is found by continuing the real root off the real axis,
which also fixes the mode's identity, node counting having no meaning for a
complex eigenfunction. Note that Eqs.~(\ref{eq:cowling}) contain $\omega$ only
as $\omega^{2}$ and so are blind to the sign convention, whereas
Eq.~(\ref{eq:cdy}) is linear in $\omega$; with the $e^{+i\omega t}$ convention
adopted here a damped mode has $\mathrm{Im}\,\omega>0$, and mixing conventions
silently reverses the sign of the damping.

As a check of the implementation, a $1.4\,M_{\odot}$ DD2 star built with both
sound speeds evaluated consistently gives $\nu_{g_{1}}=150$~Hz, within the
$130$--$270$~Hz spread of published cold nucleonic results collected in Table~I
of Ref.~\cite{Jaikumar2021}, and an $f$-mode at $2.06$~kHz. An equation of
state with a single sound speed returns the $f$-mode alone and no $g$-branch at
all, as it must.

\section{Prospects for detection}

$g$-modes couple weakly to gravitational radiation, so they are unlikely to be
seen as an isolated ringing. Two routes are more promising.

\emph{Resonant tidal excitation during inspiral.} As the orbital frequency
sweeps upwards it passes through the $g$-mode frequency, transferring orbital
energy into the mode and advancing the gravitational-wave phase~\cite{Lai1994}.
The accumulated phase error is of order $\Delta\Phi\sim0.5$--$1$~rad, large
enough to matter for waveform modelling. It is important to be honest about the
discriminating power: a hybrid star's higher $g$-mode frequency means the
resonance occurs \emph{later} in the inspiral, leaving less time to accumulate
phase, so the net $\Delta\Phi$ is comparable to that of an ordinary neutron star
even though the energy transferred is $2$--$3$ times
larger~\cite{Jaikumar2021}. Distinguishing the two therefore requires the
resonance to be localised in frequency, not merely detected.

\emph{Post-merger remnants.} A hypermassive remnant is violently excited and
hot. Here the finite-rate physics of Eq.~(\ref{eq:cdy}) is decisive: at
$T\gtrsim$ a few MeV the Urca rates reach the mode frequency, and composition
$g$-modes are damped and ultimately suppressed~\cite{Counsell2025}. Any search
must therefore target the cooler, later phase.

Third-generation detectors --- Einstein Telescope~\cite{Punturo2010} and Cosmic
Explorer~\cite{Reitze2019} --- have their best strain sensitivity in the
$100$--$1000$~Hz band where these modes lie, and the quasi-universal relations
between $f$- and $g$-mode frequencies and the tidal
deformability~\cite{ZhaoLattimer2022} offer a route to combining a marginal mode
measurement with the far better-constrained $\Lambda$. A firm $g$-mode frequency
near $500$~Hz in a star of known mass would be difficult to explain with
nucleonic matter alone.

\begin{acknowledgments}
\end{acknowledgments}

\bibliographystyle{apsrev4-2}
\bibliography{../../../docs/eos}

\end{document}
